\section{Нефункціональні вимоги}
\label{sec:nefunktsionalni-vymohy}

\subsection{Продуктивність}

\textbf{REQ-NF-001.} Сторінки системи повинні завантажуватися в прийнятний для користувача час за нормальних умов мережевого підключення. \sked{K34}

\subsection{Безпека}

\textbf{REQ-NF-002.} Паролі користувачів зберігаються виключно у захищеному вигляді~--- у вигляді хешів (наприклад, bcrypt). Доступ до захищених ресурсів вимагає активної автентифікованої сесії. \sked{K35}

\textbf{REQ-NF-007.} Авторизаційний токен зберігається в HTTP-only cookie для запобігання XSS-атакам. Всі запити, що змінюють стан системи, повинні бути захищені від CSRF-атак. \sked{K47}

\subsection{Надійність}

\textbf{REQ-NF-003.} Усі дані користувача зберігаються в базі даних. Жодна дія користувача не повинна призводити до втрати раніше введених даних без явного підтвердження видалення з боку користувача. \sked{K36}

\subsection{Зручність використання}

\textbf{REQ-NF-004.} Інтерфейс системи повинен бути адаптивним для десктопних браузерів і мобільних браузерів. Окрема оптимізація під планшет у версії~v1 не передбачена. \sked{K37}

\textbf{REQ-NF-008.} При порожньому каталозі, порожній добірці або відсутніх результатах пошуку~/ фільтру система повинна відображати інформативне повідомлення замість порожнього списку. \sked{K52}

\subsection{Сумісність}

\textbf{REQ-NF-005.} Система повинна коректно функціонувати в останніх двох версіях браузерів: Google~Chrome, Mozilla~Firefox, Microsoft~Edge та Apple~Safari (включно з мобільним Safari на iOS). \sked{K45}

\subsection{Процес розробки}

\textbf{REQ-NF-009.} Після кожного запуску тестів агент-розробник зобов'язаний записати log-файл у директорію \texttt{logs/}. Log-файл повинен містити: загальну кількість тестів, кількість passed~/ failed~/ skipped, статус кожного окремого тесту та час виконання тестового набору. \sked{K54}

\subsection{Валідація вхідних даних}

\textbf{REQ-NF-006.} При реєстрації нового користувача система повинна перевіряти:
\begin{itemize}
  \item email~--- формат відповідно до RFC~5322;
  \item пароль~--- мінімальна довжина 8~символів;
  \item нікнейм~--- від 2 до 50~символів.
\end{itemize}
При порушенні будь-якого з обмежень система виводить конкретне повідомлення про помилку, що вказує на причину відмови. \sked{K46}
